\section{The Weil form in the logarithmic variable}\label{sec:setup}

Tests are complex-valued. We use
\[
 \widehat g(z)=\int_{\R}g(x)e^{-izx}dx,\qquad
 (U_qg)(x)=g(x-q),\qquad g^*(x)=\overline{g(-x)},
\]
\[
 \xi(s)=\tfrac12s(s-1)\pi^{-s/2}\Gamma(s/2)\zeta(s),\qquad
 \Xi(z)=\xi(\tfrac12+iz).
\]
The change of variable $x=\log u$ carries $du/u$ to $dx$ and
multiplicative convolution and involution to the displayed additive
ones. We write $\Lam$ for the von Mangoldt function,
$w_n=\Lam(n)/\sqrt n$, and $Z(\Xi)$ for the zero multiset, with
multiplicity. The constant $\gamma$ below is the Euler--Mascheroni
constant.

\begin{definition}[source form]\label{def:Q}
For $g\in C_c^\infty(\R;\C)$ set
\[
 H(g)=\|g\|_2^2,\quad
 C_g(t)=\Re\int_{\R}\overline{g(x)}g(x+t)dx,\quad
 A_\pm(g)=\int_{\R}g(x)e^{\pm x/2}dx,
\]
\[
 a(t)=\frac{e^{-t/2}}{1-e^{-2t}},\quad
 c_A=\gamma+\log(8\pi)+\frac\pi2,\quad
 \D(g)=\int_0^\infty a(t)\|g(\cdot+t)-g\|_2^2dt.
\]
Define
\begin{equation}\label{eq:Q}
 \Q(g)=\D(g)-c_AH(g)
 +2\Re\bigl(A_+(g)\overline{A_-(g)}\bigr)
 -2\sum_{n\ge2}w_nC_g(\log n).
\end{equation}
The Hermitian polarisation $\B$ is antilinear in its first argument.
For use below, put
\[
 C_{g,v}(t)=\frac12\int_{\R}
  \bigl(\overline{g(x)}v(x+t)+\overline{g(x+t)}v(x)\bigr)dx;
 \qquad C_{g,g}=C_g.
\]
\end{definition}

The prime sum is finite on compact tests: if the support of a test has
diameter at most $L$, its autocorrelation is supported in $[-L,L]$.
In this convention the polarized explicit formula is
\begin{equation}\label{eq:EF}
 \B(g,v)=\sum_{z\in Z(\Xi)}
          \overline{\widehat g(\bar z)}\widehat v(z),
 \qquad g,v\in C_c^\infty(\R;\C).
\end{equation}
This is Weil's formula, in the compact-test formulation and
normalisation used in
\cite{BOMBIERI-WEILQF-2000,CCM-ZST-2025,Suzuki2023Screw}.
Indeed, $\widehat{g^**v}(z)=
\overline{\widehat g(\bar z)}\widehat v(z)$.
For an analytic check of the archimedean constant, let $\psi_{\rm d}$
denote the digamma function. Its integral representation gives
\[
 2\int_0^\infty a(t)(1-\cos(ut))dt
   =\Re\psi_{\rm d}(\tfrac14+\tfrac{iu}2)-\psi_{\rm d}(\tfrac14),
 \qquad
 \psi_{\rm d}(\tfrac14)=-\gamma-\tfrac\pi2-3\log2.
\]
Thus the Fourier multiplier of $\D-c_AH$, with Plancherel factor
$1/(2\pi)$, is
$\Re\psi_{\rm d}(1/4+iu/2)-\log\pi$.
The two pole evaluations and the two prime translations then give
\eqref{eq:Q}. This checks the normalization without testing a vector
whose form value vanishes. Smooth compact support gives absolute
convergence of the zero sum. Under RH it becomes
$\Q(g)=\sum_z|\widehat g(z)|^2$; that rewriting is not used
unconditionally. Weil's criterion is
$\Q(g)\ge0$ for every complex $g\in C_c^\infty(\R)$ if and only if RH.

\paragraph{Control space.}
Let $\Xf$ be the completion of $C_c^\infty(\R;\C)$ for
\[
 \|g\|_{\Xf}^2=\|e^{|x|}g\|_2^2+\D(g).
\]
This completion has an injective realization as a space of functions.
To see this, set
$Jg(x,t)=a(t)^{1/2}(g(x+t)-g(x))$ for $t>0$.
If $g_j$ is Cauchy in this norm, then $g_j$ converges in weighted $L^2$
and $Jg_j$ converges in $L^2(dx\,dt)$. On each strip
$\delta\le t\le T$, translation continuity in $L^2$ identifies the
second limit with $Jg$; exhaustion identifies it for all $t>0$.
In particular a zero weighted-$L^2$ limit has zero second limit.
We use this realization of the completion, not an unproved claim that
it equals every function with finite displayed norm.

The mixed translation energy is bounded by
$\D(g)^{1/2}\D(v)^{1/2}$. Moreover,
\[
 |A_\pm(g)|\le\frac2{\sqrt3}\|e^{|x|}g\|_2,
 \qquad
 |C_{g,v}(t)|\le e^{-|t|}
       \|e^{|x|}g\|_2\|e^{|x|}v\|_2.
\]
The prime coefficient sum is at most
$\sum_{n\ge2}(\log n)n^{-3/2}<5$, by comparison with the
integral of the decreasing function $(\log x)x^{-3/2}$ on $[2,\infty)$.
Consequently
\[
 |\B(g,v)|\le C_{\Xf}\|g\|_{\Xf}\|v\|_{\Xf},
 \qquad C_{\Xf}=|c_A|+14,
\]
since $1+8/3+10<14$ bounds the three other contributions.
This constructs the continuous extension of $\B$.
Translations extend boundedly, with
$\|U_qg\|_{\Xf}\le e^{|q|}\|g\|_{\Xf}$, because $\D$ is
translation-invariant. Membership of the noncompact functions used
below is established by their cutoff estimates, not assumed.

\paragraph{Finite windows.}
Write $L=\log m=2\log\lambda$. On
$J_L=[-L/2,L/2]$, the translated CCM Fourier basis is
\[
 V_n(x)=L^{-1/2}\exp(2\pi in(x+L/2)/L)\mathbf1_{J_L}(x).
\]
The full compression on $|n|\le N$ has dimension $2N+1$.
It is not a prolate eigenfunction expansion; prolate functions enter
separately as trial functions \cite{CCM-ZST-2025,ConnesMoscovici2021}.
A cell is one pair $(m,N)$, with $N$ a Fourier cutoff.
The full-window bottom is the infimum of
$\Q(g)/\|g\|_2^2$ over nonzero vectors of the window form domain.
For fixed $m$, the lowest eigenvalues of the full Fourier compressions
decrease to this infimum, denoted $\mu_\lambda$
\cite[Proposition 3.4]{CCM-ZST-2025}. A finite compression therefore
provides an upper bound, not a positivity certificate for $\mu_\lambda$.
Nonnegativity on every window is equivalent to the full Weil criterion;
the monotonicity in \cite[Corollary 3.7]{CCM-ZST-2025} also transfers
cofinal nonnegativity to all smaller windows.
Classical short-support positivity is discussed in
\cite{YOSHIDA-HERMITIAN-1992,BOMBIERI-WEILQF-2000,CC-WEILPOS-2020}.
The recent preprint \cite{XUEFENGZHU-2026} reports a larger certified
bounded-support range; those certificates are not a dependency of the
present paper and are not verified here.
